\documentclass[10pt,a4paper]{amsart}
\usepackage{amsmath,amssymb,amsthm,graphicx}
% Retained preamble for the Stacks Project chapter divisors.tex (Divisors).
% The chapter file pinned-original-source.tex carries only \input{preamble} on its first line, so no
% preamble text lives inside the chapter itself.  The Stacks Project uses one single shared preamble for
% every chapter; the snapshot stages it verbatim in sources/stacks-chow/preamble.tex.  The block below is
% copied mechanically, line for line, from that file's lines 14-18 (the Xy-pic setup) and lines 154-217
% (the native theorem declarations together with the mathematical macros).  Nothing was added, removed or
% redefined.  The wrapper class and the amsmath/amssymb/amsthm/graphicx base packages are supplied by the
% assembler, so the class/\externaldocument part of the shared preamble is not repeated here.
% For commutative diagrams we use Xy-pic
\usepackage[all]{xy}

% We use 2cell for 2-commutative diagrams.
\xyoption{2cell}
% Theorem environments.
%
\theoremstyle{plain}
\newtheorem{theorem}[subsection]{Theorem}
\newtheorem{proposition}[subsection]{Proposition}
\newtheorem{lemma}[subsection]{Lemma}

\theoremstyle{definition}
\newtheorem{definition}[subsection]{Definition}
\newtheorem{example}[subsection]{Example}
\newtheorem{exercise}[subsection]{Exercise}
\newtheorem{situation}[subsection]{Situation}

\theoremstyle{remark}
\newtheorem{remark}[subsection]{Remark}
\newtheorem{remarks}[subsection]{Remarks}

\numberwithin{equation}{subsection}

% Macros
%
\def\lim{\mathop{\mathrm{lim}}\nolimits}
\def\colim{\mathop{\mathrm{colim}}\nolimits}
\def\Spec{\mathop{\mathrm{Spec}}}
\def\Hom{\mathop{\mathrm{Hom}}\nolimits}
\def\Ext{\mathop{\mathrm{Ext}}\nolimits}
\def\SheafHom{\mathop{\mathcal{H}\!\mathit{om}}\nolimits}
\def\SheafExt{\mathop{\mathcal{E}\!\mathit{xt}}\nolimits}
\def\Sch{\mathit{Sch}}
\def\Mor{\mathop{\mathrm{Mor}}\nolimits}
\def\Ob{\mathop{\mathrm{Ob}}\nolimits}
\def\Sh{\mathop{\mathit{Sh}}\nolimits}
\def\NL{\mathop{N\!L}\nolimits}
\def\CH{\mathop{\mathrm{CH}}\nolimits}
\def\proetale{{pro\text{-}\acute{e}tale}}
\def\etale{{\acute{e}tale}}
\def\QCoh{\mathit{QCoh}}
\def\Ker{\mathop{\mathrm{Ker}}}
\def\Im{\mathop{\mathrm{Im}}}
\def\Coker{\mathop{\mathrm{Coker}}}
\def\Coim{\mathop{\mathrm{Coim}}}

% Boxtimes
%
\DeclareMathSymbol{\boxtimes}{\mathbin}{AMSa}{"02}

%
% Macros for moduli stacks/spaces
%
\def\QCohstack{\mathcal{QC}\!\mathit{oh}}
\def\Cohstack{\mathcal{C}\!\mathit{oh}}
\def\Spacesstack{\mathcal{S}\!\mathit{paces}}
\def\Quotfunctor{\mathrm{Quot}}
\def\Hilbfunctor{\mathrm{Hilb}}
\def\Curvesstack{\mathcal{C}\!\mathit{urves}}
\def\Polarizedstack{\mathcal{P}\!\mathit{olarized}}
\def\Complexesstack{\mathcal{C}\!\mathit{omplexes}}
% \Pic is the operator that assigns to X its picard group, usage \Pic(X)
% \Picardstack_{X/B} denotes the Picard stack of X over B
% \Picardfunctor_{X/B} denotes the Picard functor of X over B
\def\Pic{\mathop{\mathrm{Pic}}\nolimits}
\def\Picardstack{\mathcal{P}\!\mathit{ic}}
\def\Picardfunctor{\mathrm{Pic}}
\def\Deformationcategory{\mathcal{D}\!\mathit{ef}}

% --- appended 2026-10-12 (Chief): the project own definitions of the reference/slogan/history
% display environments, copied mechanically from the shared Stacks preamble lines 8-12
% (\usepackage{verbatim} plus the three \newenvironment lines, which the chapter preamble keeps
% verbatim from the project preamble). Nothing added, removed or redefined: these environments
% carry no mathematical content and the project itself typesets them as comments. Without these
% three lines a counted statement whose own text uses them fails with Environment slogan undefined.
\usepackage{verbatim}
\newenvironment{reference}{\comment}{\endcomment}
\newenvironment{slogan}{\comment}{\endcomment}
\newenvironment{history}{\comment}{\endcomment}
\usepackage{cleveref}
\crefname{originallemma}{original lemma}{original lemmas}
\Crefname{originallemma}{Original lemma}{Original lemmas}
\makeatletter
\newlabel{lemma-open-subscheme-UFD}{{0BDA}{source}{}{sourcepointer-1}{}}
\newlabel{lemma-open-subscheme-UFD@cref}{{[originallemma][0][]0BDA}{[1][0][]source}{}{}{}}
\makeatother
\title{The Picard group of projective space over a unique factorisation domain is the integers, generated by the twisting sheaf}
\author{The Stacks Project Authors}
\date{stacks-project:divisors.tex:local-snapshot-20261012; GNU FDL 1.2 or later}
\begin{document}
\maketitle
\noindent\textit{Editorial note.} All delivered mathematical text is the author's own, copied from the pinned Stacks Project chapter file pinned-original-source.tex; no statement or proof step was rewritten, reordered or removed.\par\medskip
\noindent\textit{Editorial note.} Following the multi-claim convention for this source, source\_url is the source-level URL https://stacks.math.columbia.edu/ and the source label lemma-Pic-projective-space-UFD is carried in claim\_key and here rather than in a per-claim tag URL.\par\medskip
\noindent\textit{Editorial note.} The preamble is the chapter's shared Stacks Project preamble, retained as sources/stacks-divisors/retained-preamble.tex: a line-for-line excerpt of the shared Stacks preamble staged with this snapshot. No macro was invented or redefined, and this card adds no preamble text of its own.\par\medskip
\noindent\textit{Editorial note.} Statement span 7254-7264 and proof span 7266-7301 come from the halfA pool card operations/pool-cards/stacks-divisors-halfA.json and were re-read in the pinned chapter file before this card was written.\par\medskip
\noindent\textit{Editorial note.} The proof works with the author's notation for the standard charts, the units, and the exponents of the coordinate ratios; the displays are preserved unchanged.\par\medskip
\noindent\textit{Editorial note.} 1 cross-references occur inside the counted spans; they are listed in external\_references as navigation pointers only.\par\medskip
\noindent\textit{Allowed standard prerequisites (editorial).}\begin{itemize}
\item The Picard group of an open subscheme of a unique factorisation domain's affine space is trivial, cited from the same chapter. Cited at proof lines 7266-7268. Shows that the standard affine charts of projective space have trivial Picard group, so invertible modules are given by cocycles for that cover.
\item The description of the units of a localization as the product of the units of the base with the powers of the localizing element, and the cocycle description of invertible modules with respect to a cover. Used at proof lines 7272-7288. Turns the cocycle into a pair of data, a unit-valued part and an integral exponent part, for each overlapping pair of charts.
\item The trivializing section of a twist of the structure sheaf on a standard chart and the corresponding cocycle. Used at proof lines 7295-7301. Identifies the twisting sheaf with the class of a cocycle consisting of powers of the ratios of coordinates, matching the classification obtained.
\end{itemize}\medskip
\ifdefined\hypertarget\hypertarget{sourcepointer-1}{}\fi
\noindent\textit{Source reference.} Original lemma 0BDA: focused-seeds/AGP/sources/stacks-snapshot-20261011/tags.txt, source lines 14308--14308. Same-chapter reference used in the retained author text at source line(s) 7267. It is supplied only as a navigation pointer so that the retained text compiles; the counted item delivers the author's own statement and the author's complete proof, and no part of the referenced item is copied in or substituted for the counted argument.\par\medskip
\begin{lemma}
\label{lemma-Pic-projective-space-UFD}
Let $R$ be a UFD. The Picard group of $\mathbf{P}^n_R$
is $\mathbf{Z}$. More precisely, there is an isomorphism
$$
\mathbf{Z} \longrightarrow \Pic(\mathbf{P}^n_R),\quad
m \longmapsto \mathcal{O}_{\mathbf{P}^n_R}(m)
$$
In particular, the Picard group of $\mathbf{P}^n_k$ of projective
space over a field $k$ is $\mathbf{Z}$.
\end{lemma}

\begin{proof}
By Lemma \ref{lemma-open-subscheme-UFD} the Picard groups of the opens
$D_+(T_i) \cong \mathbf{A}^n_R$ are trivial.
Thus if $\mathcal{L}$ is an invertible module on
$\mathbf{P}^n_R$, then it is given by a $1$-cocycle with values in
the sheaf of invertible functors $\mathcal{O}^*$ for the open covering
$$
\mathbf{P}^n_R = \bigcup\nolimits_{i = 0, \ldots, n} D_+(T_i)
$$
Observe that for $i \not = j$ we have
$$
\mathcal{O}^*(D_+(T_i) \cap D_+(T_j)) =
\mathcal{O}^*(D_+(T_iT_j)) = \left(R[T_0, \ldots, T_n]_{(T_iT_j)}\right)^* =
R^* \times (T_i/T_j)^\mathbf{Z}
$$
Thus such a cocycle $(g_{ij})$ is given by units
$$
g_{ij} = u_{ij} (T_i/T_j)^{e_{ij}}
$$
with $u_{ij} \in R^*$ and $e_{ij} \in \mathbf{Z}$
satisfying the cocycle condition. The cocycle condition
over $D_+(T_iT_jT_k)$ for $\#\{i, j, k\} = 3$ tell us that
$$
u_{ik} = u_{ij} u_{jk} \quad\text{and}\quad
e_{ik} = e_{ij} = e_{jk}
$$
Whence $u_{ij} = u_{i1} u_{j1}^{-1}$ is a boundary.
Thus all isomorphism classes of invertible modules
are given by taking the cocycle with
$$
g_{ij} = (T_i/T_j)^e
$$
for some $e \in \mathbf{Z}$. Since $\mathcal{O}(n)$ has trivializing
section $T_i^n$ over $D_+(T_i)$ we see that the corresponding
cocycle of $\mathcal{O}(n)$ is $(T_i/T_j)^n$ and the proof is complete.
\end{proof}

\end{document}
